% ---------------------------------------------------------------
\section{Mathematical framework: dual spectral ranking}
\label{sec:networks}
% ---------------------------------------------------------------

We begin by selecting $N_S$ sources $\mathfrak{S}$ in a set of fields, disciplines or categories $\mathfrak{F}$.
Then we select $N_w$ works $\mathfrak{W}$ published in $\mathfrak{S}$ within the census
window years $t_c = [t^c_0, t^c_1]$ or target years $t_t = [t^t_0, t^t_1]$.
From the reference lists of the set of census works $\mathfrak{W}^c$ we identify the set
of target works $\mathfrak{W}^t$ and form the filtered reference set
$\mathfrak{R}=\{r(i,j):w_i\in\mathfrak{W}^c,\,w_j\in\mathfrak{W}^t\}$.
Author affiliations in $\mathfrak{W}$ establish a set of institutions $\mathfrak{I}'$ and
authors $\mathfrak{A}'$. In practice many institutions and some sources publishing only
a few works per year can be dropped from the corpus. Let $\bar{w}_I$ be the mean annual
census-window work count for institution $x$ and $\bar{w}_S$ the corresponding work count for source $s$.
We determine lower limits $\tau_I$ and $\tau_S$ by examining the respective frequency distributions and define
$\mathfrak{I}\subseteq\mathfrak{I}'$ as the institutions with $\bar{w}_I\geq\tau_I$,
and $\mathfrak{S}\subseteq\mathfrak{S}'$ as the sources with $\bar{w}_S\geq\tau_S$,
with corresponding author set $\mathfrak{A}\subseteq\mathfrak{A}'$.
Works that lose all their authorships under the institution filter, or whose source is
dropped by the source filter, are removed from the corpus. The choice to apply $\tau_I$ and $\tau_S$ once across
the corpus or else within each time window is examined below.

Let $R_i=|\{j:(i, j)\in\mathfrak{R}\}|$ be the number of references from work
$w_i$ that land in the filtered set $\mathfrak{R}$, and
$\bar{R}=|\mathfrak{R}|/N_w$ the mean of $R_i$ over $\mathfrak{W}$. Since filtered reference lists
have different lengths, we consider two cases for weighting references,

\begin{equation}
	\rho_i \;=\;
	\begin{cases}
		 1 & \text{(full reference counting)} \\
		\bar{R}/R_i & \text{(fractional reference counting).}
	\end{cases}
	\label{eq:rho}
\end{equation}
Under full reference counting a work with twice as many filtered references has twice
the prominence; under fractional reference counting every work's filtered references count
equally \citep{zittModifyingJournalImpact2008}.

A work has a retained source and one or more institutions. For journal spectral ranking, it is straightforward to link the work to its journal. For institution spectral ranking it is necessary to consider institution fractionation, while in dual spectral ranking, both institution fractionation and the fractionation between the journal and the institutions must be considered. 

There are two models for fractionating a work's institutions (i.e. the \textit{institution weight})
\begin{equation}
  \omega_{ix} \;=\;
  \begin{cases}
  \displaystyle\sum_{\ell:\,x\in\mathcal{I}_{i\ell}}
  \frac{1}{a_i}\,\frac{1}{x_{i\ell}}, & \text{(author-fractional)} \\[10pt]
  \dfrac{1}{n_i}, & \text{(institution-fractional).}
  \end{cases}
  \label{eq:omega}
\end{equation}
Here, $\mathcal{I}_{i\ell}$ is the set of institutions of author $\ell$ of work $w_i$, $x_{i\ell}=|\mathcal{I}_{i\ell}|$ is its cardinality, $a_i$ is the author count of work $w_i$, and $n_i = |\bigcup_\ell \mathcal{I}_{i\ell}|$ is the number of distinct institutions of work $w_i$. Authorships with no institution are automatically excluded, so $x_{i\ell}\ge 1$. Under author-fractional counting $\sum_x\omega_{ix}=1$ \citep{waltmanFieldnormalizedCitationImpact2015}. Institution pruning on annual output $\tau_I$ precedes fractional counting; this choice boosts a work's remaining institutions but diminishes the work's total contribution when some institutions are dropped.

In dual ranking, the parameter $\chi$ fractionates between institutions $(\chi)$ and source $(1-\chi)$. We consider cases $\chi = 0.5$ (unit balance) and $\chi=N_I/(N_S+N_I) = \chi^*$ (unit count balance).

We construct the set of \emph{units} [sources or institutions;
\cite{pinskiCitationInfluenceJournal1976}]
$\mathfrak{P}=\mathfrak{S}\ \cup\ \mathfrak{I}$ with $N=N_S+N_I$ elements,
ordered so that sources occupy indices $1,\dots,N_S$ and institutions
indices $N_S+1,\dots,N$. Each reference $r(i,j) \in \mathfrak{R}$ connects units
associated with $w_i$ to units associated with $w_j$,
across four unit-type pairs: source--source ($SS$), source--institution
($SI$), institution--source ($IS$), and institution--institution ($II$).

For each work $w_i$ we define the source membership one-hot vector
$\mathbf{b}_S(i)\in\mathbb{R}^{N_S}$ with $[\mathbf{b}_S(i)]_x=\mathbf{1}[w_i\in x]$,
and the institution weight distribution vector
$\mathbf{b}_I(i)\in\mathbb{R}^{N_I}$ with $[\mathbf{b}_I(i)]_y=\omega_{iy}$
from \eqref{eq:omega}. The vectors $\mathbf{b}$ yield the projection across unit
pairs, accumulated as four block matrices over $\mathfrak{R}$ via
rank-1 outer products,
\begin{equation}
	\mathbf{C}_{XY} \;=\; \sum_{(i,j)\,\in\,\mathfrak{R}}
	\rho_i\;\mathbf{b}_X(i)\,\mathbf{b}_Y(j)^\top,
	\qquad XY\in\{SS,\,SI,\,IS,\,II\}.
	\label{eq:Cblock}
\end{equation}
Here, $\mathbf{b}_X(i)$ distributes the referencing side across
units of type $X$, and $\mathbf{b}_Y(j)$ distributes the referenced side
across units of type $Y$. Each entry $[\mathbf{C}_{XY}]_{pq}$ is therefore proportional
to the sum over $\mathfrak{R}$ of $\rho_i$-weighted fractional references from unit $p$ of type $X$ to unit $q$ of type $Y$. Note that throughout this paper, rows of $\mathbf{C}$ index the referencing side and columns index the referenced side. Row sums are therefore total references \emph{given} and column sums are total references \emph{received}.

The reference matrix assembles the four blocks:
\begin{equation}
	\mathbf{C}(\mathbf{m},\rho,\chi) \;=\;
	\begin{bmatrix}
		m_{SS}\,\mathbf{C}_{SS} &
		m_{SI}\,\chi(1-\chi)\,\mathbf{C}_{SI} \\[4pt]
		m_{IS}\,\chi(1-\chi)\,\mathbf{C}_{IS} &
		m_{II}\,\mathbf{C}_{II}
	\end{bmatrix}.
	\label{eq:C}
\end{equation}
The binary \textit{block mask} $\mathbf{m}=(m_{SS},m_{SI},m_{IS},m_{II})\in\{0,1\}^4$
is bibliometrically meaningful for source-only $(1000)$, institution-only $(0001)$,
bipartite $(0110)$, and joint-full $(1111)$ combined blocks.
The institution-source fractionation $\chi$ weights the bipartite and joint masks.
As a consequence of filtering $\mathfrak{R}$ to census-to-target references, the blocks $\mathbf{C}_{XY}$ exclude (a) census-work references not giving
attention to target works, (b) non-census works giving attention to
target works, and (c) authorships that include an omitted institution.

\paragraph{Ordinal spectral ranking.}

Let $r_p=\sum_q C_{pq}$ be the sum of fractionated references given by unit $p$
and $\mathbf{D}_r=\mathrm{diag}(r_p)$ the corresponding diagonal matrix.
Row-normalising $\mathbf{C}$ gives the row-stochastic matrix
$\mathbf{H}=\mathbf{D}_r^{-1}\mathbf{C}$,
ensuring that, when $\rho_i = 1$, differences in reference intensity across units are
removed. They are reintroduced later
via the influence scaling \eqref{eq:vstar}.
Dangling units (units with $r_p=0$) and units with no incoming references are iteratively excluded from the strongly connected component on which ranking is computed, as described below.

The masks source-only $(1000)$, institution-only $(0001)$, and full-joint $(1111)$
each yield a square $\mathbf{H}$ with the spectral
ranking $\boldsymbol{\pi}^*$ \citep[see][for an overview]{vignaSpectralRanking2016}
\begin{equation}
	(\mathbf{I} - \mathbf{H}^\top)\boldsymbol{\pi}^* = \mathbf{0},
	\qquad
	\mathbf{1}^\top\boldsymbol{\pi}^*=1,
	\quad
	\boldsymbol{\pi}^*\ge 0.
	\label{eq:eig}
\end{equation}
If $\mathbf{H}$ is primitive (irreducible and aperiodic) the ranking is unique by the
Perron--Frobenius theorem. Our numerical implementation tests primitivity
for every $\mathbf{H}$, pre-processed as explained below, and has not encountered an imprimitive case.
\cite{gellerCitationInfluenceMethodology1978} showed that
$\boldsymbol{\pi}^*$ is the stationary distribution of the
Markov chain with transition matrix $\mathbf{H}$, and that it coincides
with the \textit{influence weight} of \cite{pinskiCitationInfluenceJournal1976}
scaled by the reference-count proportion $r_p / \sum_q r_q$.

\cite{katzNewStatusIndex1953} proposed attenuating walks (chains of reference-matrix weights
$c_{ab}\rightarrow c_{bc}\rightarrow c_{cd}$) of length $k$ by a factor $\alpha^k$, $\alpha\in(0,1)$, so that distant
links count for less than proximate ones.
Summing walks of all lengths and endowing every unit with a
uniform initial score of $1$ yields the Katz resolvent equation
\[
	(\mathbf{I} - \alpha\mathbf{H}^\top)\boldsymbol{\pi} = \mathbf{1}.
\]
Because $\|\mathbf{H}\|_\infty\leq 1$, the spectral radius $\rho(\alpha\mathbf{H}^\top)<1$ for any $\alpha\in(0,1)$, so $(\mathbf{I}-\alpha\mathbf{H}^\top)$ is invertible by Neumann series.

\cite{hubbellInputoutputApproachClique1965} generalised this approach by replacing
the uniform endowment $\mathbf{1}$ with an informative prior
$\boldsymbol{\mu}$ over units, representing exogenous information about
standing independent of reference structure,
\[
	(\mathbf{I} - \alpha\mathbf{H}^\top)\boldsymbol{\pi} = \boldsymbol{\mu}.
\]
We explore two priors, normalised so that $\mathbf{1}^\top\boldsymbol{\mu}=1$:
$\mu_p = 1/N$ for every unit $p\in\mathfrak{P}$ (uniform prior) and $\mu_p=1/(2N_S)$ for sources and $\mu_p=1/(2N_I)$ for institutions (unit-weighted prior).
Note that if $\alpha < 1$ and $\mu=0$ the resolvent $(\mathbf{I}-\alpha\mathbf{H}^\top)\boldsymbol{\pi}=\mathbf{0}$ admits only the trivial solution $\boldsymbol{\pi}=\mathbf{0}$ and is not useful as a ranking.

Combining these yields the general (Katz--Hubbell) resolvent equation
\begin{equation}
	(\mathbf{I} - \alpha\mathbf{H}^\top)\boldsymbol{\pi}
	\;=\;
	(1-\alpha)\boldsymbol{\mu},
	\qquad
	\mathbf{1}^\top\boldsymbol{\pi}=1,
	\quad
	\boldsymbol{\pi}\ge 0,
	\label{eq:spectral}
\end{equation}
which recovers $\boldsymbol{\pi}^*$ in the limit $\alpha\to 1$ when
$\mathbf{H}$ is primitive.
The prior $\boldsymbol{\mu}$ has a substantive interpretation as exogenous standing independent of reference structure. Together with $\alpha<1$ it additionally guarantees a unique nonnegative solution to \eqref{eq:spectral} even when $\mathbf{H}$ is imprimitive: $\alpha<1$ makes $(\mathbf{I}-\alpha\mathbf{H}^\top)$ invertible, and a nonzero $\boldsymbol{\mu}$ lifts the equation off the homogeneous system. If $\mathbf{C}$ were a binary adjacency matrix with the flat prior, this would be the PageRank spectral ranking \citep{brinAnatomyLargeScaleHypertextual1998}.

The spectral ranking $\boldsymbol{\pi}$ is the steady-state probability vector of the Markov matrix $\mathbf{H}$; we refer to $\pi_p$ as the \emph{ranking probability} of unit $p$.
The \emph{influence} $v_p$ is computed as \citep{palacios-huertaMeasurementIntellectualInfluence2004},
\begin{equation}
	\mathbf{v}
	\;=\;
	A\,\mathbf{A}^{-1}\boldsymbol{\pi},
	\qquad
	\mathbf{A}=\mathrm{diag}(a_p),
	\qquad
	A = \mathbf{a}^{\top}\mathbf{1} = \textstyle\sum_p a_p,
	\label{eq:vstar}
\end{equation}
i.e.\ $v_p = (A/a_p)\,\pi_p$. 
Here, the unit work count $a_p = |\{i:w_i\in p\}|$ where $p$ is a specific source or institution,
and $\boldsymbol{\pi}$ denotes the solution to \eqref{eq:spectral}.
The scalar $A$ normalises $\mathbf{v}$ so that its $a_p$-weighted mean equals one:
$\mathbf{a}^{\top}\mathbf{v}/A=\mathbf{1}^{\top}\boldsymbol{\pi}=1$.
A unit with $v_p=1$ has the corpus-average influence, and
$v_p=2$ indicates twice the corpus-average influence.  The quantity $v$ is the dual ranking counterpart of influence per publication \citep{pinski} and article influence score \cite{bergstromEigenfactorMetricsFigure2008}.
\cite{palacios-huertaMeasurementIntellectualInfluence2004} show that
$\mathbf{v}$ is uniquely characterised by making several reasonable axiomatic
choices, although \cite{serranoMeasurementIntellectualInfluence2004} offers an impossibility result by proposing a reasonable additional axiom that breaks the uniqueness.

\paragraph{Bipartite ranking.}
With $\mathbf{m}=(0110)$, equation~\eqref{eq:C} yields block off-diagonal matrices $\mathbf{C}$ (and $\mathbf{H}$) forming a directed bipartite network between sources~($S$) and institutions~($I$).
This mask characterises our dual ranking problem and is used as the baseline model in Section~\ref{sec:results}.
Partitioning rows and columns into the $S$ and $I$ blocks, $\mathbf{H}$ takes the form
\begin{equation}
  \mathbf{H}
  \;=\;
  \begin{pmatrix} \mathbf{0} & \mathbf{H}_{SI} \\ \mathbf{H}_{IS} & \mathbf{0} \end{pmatrix},
  \label{eq:H_bipartite}
\end{equation}
where $\mathbf{H}_{SI}$ connects sources to institutions and $\mathbf{H}_{IS}$ does the reverse.
Because every path alternates $S\to I\to S\to\cdots$ this structure makes $\mathbf{H}$ periodic with period~2, and thus imprimitive.
Squaring $\mathbf{H}$ collapses two-step paths to separate unit types, giving one-unit relations
\[
  \mathbf{H}^2
  \;=\;
  \begin{pmatrix} \mathbf{M}_S & \mathbf{0} \\ \mathbf{0} & \mathbf{M}_I \end{pmatrix},
  \qquad
  \mathbf{M}_S = \mathbf{H}_{SI}\mathbf{H}_{IS},
  \quad
  \mathbf{M}_I = \mathbf{H}_{IS}\mathbf{H}_{SI}.
\]
Both $\mathbf{M}_S$ and $\mathbf{M}_I$ are row semi-stochastic products of row semi-stochastic matrices.

We adopt the convention that a single $SI$ or $IS$ step is damped by $\sqrt{\alpha}$, so the round-trip $S\to I\to S$ incurs attenuation $\sqrt{\alpha}\cdot\sqrt{\alpha}=\alpha$, matching the one-step attenuation in the unipartite modes. This ensures the spectral radius of $\alpha\mathbf{M}_S$ equals that of the corresponding unipartite operator and makes $\alpha$ comparable across masks.
Applying this substitution to \eqref{eq:spectral},
partitioning $\boldsymbol{\pi}=(\boldsymbol{\pi}_S,\boldsymbol{\pi}_I)$ and $\boldsymbol{\mu}=(\boldsymbol{\mu}_S,\boldsymbol{\mu}_I)$,
and eliminating $\boldsymbol{\pi}_I$ yields the bipartite resolvent for sources,
\begin{equation}
  \bigl(\mathbf{I}-\alpha\mathbf{M}_S^\top\bigr)\boldsymbol{\pi}_S
  \;=\;
  (1-\sqrt{\alpha})\bigl(\boldsymbol{\mu}_S+\sqrt{\alpha}\,\mathbf{H}_{IS}^\top\boldsymbol{\mu}_I\bigr).
  \label{eq:kh_source}
\end{equation}
When $\alpha<1$ and $\boldsymbol{\mu}\ne\mathbf{0}$, the right-hand side links institutions to sources via $\mathbf{H}_{IS}^\top$. The institution ranking is the matrix--vector product
\begin{equation}
  \boldsymbol{\pi}_I
  \;=\;
  \sqrt{\alpha}\,\mathbf{H}_{SI}^\top\boldsymbol{\pi}_S+(1-\sqrt{\alpha})\boldsymbol{\mu}_I.
  \label{eq:kh_inst}
\end{equation}

\paragraph{Computation.}
Defining the modified matrix
\begin{equation}
	\mathbf{\tilde{H}}
	\;=\;
	\alpha\mathbf{H}+(1-\alpha)\boldsymbol{\mu}\mathbf{1}^\top,
	\label{eq:Hstar}
\end{equation}
the resolvent
$\boldsymbol{\pi}=\mathbf{\tilde{H}}^{\top}\boldsymbol{\pi}$
is equivalent to \eqref{eq:spectral}. This is the general form of
our numerical problem, though most results reported below use $\alpha=1$ and $\boldsymbol{\mu}=\mathbf{0}$, so $\mathbf{\tilde{H}}=\mathbf{H}$.

We construct many realisations of $\mathbf{\tilde{H}}$ to test unit pairings, subfield restrictions, other parameter choices, and for bootstrap uncertainty estimation. Each realisation entails a unit--unit reference network with one giant strongly connected component (SCC) together with a few, typically singleton, components that are disjoint from the main network. We retain the SCC, and drop the disconnected components and any units whose row sum or column sum on the SCC is zero, alternating between component, row and column drops until the residual $\mathbf{\tilde{H}}$ is strictly row-stochastic. Primitivity is then verified numerically as described below.

$\mathbf{\tilde{H}}$ is a sparse matrix constructed from the filtered and weighted edge list $\mathcal{E}$ accumulated over $\mathfrak{R}$, stored in compressed sparse row (CSR) format. We use two numerical approaches, applied to different regimes:
\begin{itemize}
	\setlength{\itemsep}{0pt}
	\item For runs with $\alpha=1$, $\boldsymbol{\mu}=\mathbf{0}$ including the baseline and bootstrap, the dominant eigenpair of $\mathbf{\tilde{H}}^\top$ is computed with \texttt{scipy.sparse.linalg.eigs}, which also returns the next few eigenvalues and thereby verifies primitivity (a unique eigenvalue of modulus~1).
	\item For $\alpha<1$ sensitivity runs, the power iteration
\begin{equation}
	\boldsymbol{\pi}^{(k+1)}
	\;=\;
	\mathbf{\tilde{H}}^{\top}\boldsymbol{\pi}^{(k)}
	\label{eq:power}
\end{equation}
is used. Because $(1-\alpha)\boldsymbol{\mu}\mathbf{1}^\top$ makes $\mathbf{\tilde{H}}$ primitive, power iteration converges to $\boldsymbol{\pi}$ from any strictly positive starting vector $\boldsymbol{\pi}^{(0)}$.
\end{itemize}
In either approach, the influence scaling \eqref{eq:vstar} is applied after $\boldsymbol{\pi}$ has been computed.
